\sectionguard
\section{The Image and What Its Examinations Establish}

\subsection{The rule this section follows}

More untruth circulates about the Guadalupe image than about any
other object of Catholic devotion, and almost all of it circulates in
good faith. The rule adopted here is simple and is applied without
exception: a claim about the physical image is stated only if this
study reached a source for it, the source is named, and the level at
which the source stands is stated. Where the popular claim could not
be traced to any reachable source, that is recorded as a negative
finding rather than passed over. Nothing is asserted about the
image's origin, and no conclusion about the events of 1531 is drawn
from any material fact.

Two prior points bound the whole discussion. First, the image is a
venerated object continuously in cult, framed and glazed for most of
its documented history, and it has been physically examined openly
on very few occasions. Second, the Church has never made any material
property of the image an object of any act of teaching or of judgment
whatever. No papal, conciliar, or dicasterial text known to this
study asserts that the image is of non-human origin.

\subsection{The documented examinations}

\subsubsection*{1666: the painters, and the reverse}

The 1666 \work{Informaciones jur\'idicas} (\S2.2) included opinions
of painters, taken with the glass case opened. On that occasion the
back of the cloth was seen---one of only three occasions on which
this is known to have happened. Francisco de Florencia, S.J., who
became the tradition's great seventeenth-century expositor, later
recorded what he and the canon Francisco de Siles saw there: not a
shadow image showing through the thin cloth, but blotches of colour
that seemed to them like the pressed juice of various flowers and
leaves---dark green of lily leaf, snowy white, purple of iris, pink
of rose, blue of violet, yellow of broom---distinct and unconfused,
at which they marvelled for some
time.\endnote{Jaime Cuadriello, ``Cifra, signo y artilugio: el
`ocho' de Guadalupe,'' \work{Anales del Instituto de Investigaciones
Est\'eticas} 39, no.~110 (2017): 155--206,
\url{https://doi.org/10.22201/iie.18703062e.2017.1.2593}, read
2026-07-25 in the SciELO M\'exico open-access full text at
\url{https://www.scielo.org.mx/scielo.php?script=sci_arttext&pid=S0185-12762017000100155},
where Florencia's passage is quoted at length and Cuadriello observes
that ``tan s\'olo en 1666 y 1751, que se sepa, alguien se pudo asomar
al reverso.'' Florencia's text was not consulted directly. The
paraphrase above renders Florencia's list; it is not offered as a
translation of his wording.}

That testimony is doubly instructive. It is an apparitionist eyewitness
saying that the reverse shows coloured staining that had soaked
through, and it is the origin of the durable belief that the image was
imprinted by the juice of flowers.

\subsubsection*{1751--1756: Miguel Cabrera and six colleagues}

On 30 April 1751 the painter Miguel Cabrera, with six colleagues,
examined the image at length with the glass removed, confirmed its
measurements, made colour tests (\term{templas}) directly on the
paint, and studied the technique and materials. He published his
opinion five years later as \work{Maravilla americana} (1756), and it
was the report sent to Rome in support of the petition for a proper
Office and Mass. Its practical consequence is documented and
unambiguous: after 1751 the colouring of the copies produced by
Cabrera and his contemporaries changed markedly, in flesh tones and
in dress, because the painters had at last seen the original
closely.\endnote{Cuadriello, ``Cifra, signo y artilugio,'' at the
opening of the article (same witness). Cabrera's book is Miguel
Cabrera, \work{Maravilla americana, y conjunto de raras maravillas,
observadas con la direccion de las reglas de el arte de la pintura en
la prodigiosa imagen de Nuestra Sra.\ de Guadalupe de Mexico}
(Mexico: imprenta del Real y m\'as antiguo Colegio de San Ildefonso,
1756); the Getty Research Institute copy is digitized at
\url{https://archive.org/details/maravillaamerica00cabr} (metadata
read 2026-07-25). Cabrera's own narrative of the 30 April 1751
examination---that the chapter convened the most credited painters in
Mexico, that they were made to observe the image at leisure without
the glass, and that each was to judge by the rules of his art whether
such marvels could be the work of human industry---is reprinted in
\work{Colecci\'on de obras y opusculos} (Madrid, 1785),
\url{https://archive.org/details/coleccindeobrasy00unkn}, read
2026-07-25.}

Cabrera also reported a feature that became famous: what he took to
be a numeral eight near the right foot, which he read as pointing
either to the octave of the Immaculate Conception or to the image as
the eighth wonder of the world.

\subsubsection*{1787: Bartolache and the Academy painters}

In 1787 the physician and mathematician Jos\'e Ignacio Bartolache
conducted the most methodologically ambitious examination before the
twentieth century. With the abbot's permission he made three
inspections with the case open, accompanied by a notary and by five
qualified painters connected with the Real Academia de San Carlos.
He addressed the question of the support material, arguing that the
cloth was not maguey (\term{ixtle}) but the fibre of a palm
(\term{iczotl}); he had ayates woven under his own supervision by
Indigenous weavers and had copies painted on them between 6 February
and 14 March 1787 as a controlled comparison; and he had his painters
examine Cabrera's ``eight,'' who reported in a notarial act of
1~March 1787 that it was nothing special and copied it identically.
He published the results posthumously in 1790.\endnote{Cuadriello,
``Cifra, signo y artilugio'' (same witness), for the three
inspections, the five Academy painters, the notarial act of 1 March
1787 and its verdict on the ``eight'' (``que no es cosa especial, y
le copiaron id\'entico''), the \term{iczotl} thesis, and the
existence of an \term{imprimatura} beneath the paint surface.
Bartolache's book is \work{Manifiesto satisfactorio anunciado en la
Gazeta de Mexico (Tom.~I.\ N\'um.~53.): Opusculo Guadalupano}
(Mexico: Felipe de Z\'u\~niga y Ontiveros, 1790); the John Carter
Brown Library copy is digitized at
\url{https://archive.org/details/manifiestosatisf00bart} (metadata
read 2026-07-25). Bartolache's book was not read for this study; his
programme is reported from Cuadriello.}

\subsubsection*{1887--1895: the crown}

The image's material history includes one episode of change that is
documented beyond dispute and is worth recording precisely because it
is so seldom mentioned. The image had long borne a painted crown. In
January 1887, while photographs were being taken with the glass open
in preparation for the canonical coronation, it was noticed that the
crown was gone. Accusations flew in the Mexican press and among the
clergy; the promoter of the coronation swore that neither he nor the
archbishop had had any part in it; others argued publicly that its
disappearance was itself a miracle by which the Virgin signified her
willingness to be crowned anew. On 30 September 1895, when the image
was returned to the collegiate church for the coronation, an
instrument was drawn up before many witnesses recording that no crown
existed on it and there was no trace that one ever had.\endnote{Edmundo
O'Gorman, \work{Destierro de sombras}, Appendix~VII, nos.~23--24,
27--28, 32, 62 and 64, read 2026-07-25 at
\url{https://historicas.unam.mx/publicaciones/publicadigital/libros/222c/222c_04_05_07_ApendiceSeptimo.pdf}.
O'Gorman's sources there are the letter of Antonio Plancarte y
Labastida to Bishop Carrillo y Ancona of 19 June 1887 (as printed by
Jos\'e Bravo Ugarte), Gabino Ch\'avez, \work{Celeste y Terrestre o las
dos coronas guadalupanas} (1895), and the \work{Calendario del m\'as
antiguo Galv\'an} for 1895. O'Gorman states that the crown ``fue
clandestinamente borrada'' and reports the suspicion that fell on the
coronation's promoters; the identity of whoever removed it is not
established and is not asserted here.}

Nothing follows from this about the image's origin. What follows is
methodological: the object now venerated is not, in every visible
detail, identical to the object venerated in 1750, and any argument
that treats the present surface as untouched since 1531 must reckon
with that.

\subsubsection*{1982: the conservation examination}

On the night of 27 September 1982 the restorer Jos\'e Sol Rosales
carried out a consolidation and surface cleaning at the initiative of
Jorge Guadarrama, conservator of the image and director of the
Basilica museum, in the presence of the then abbot Guillermo
Schulenburg and his archpriest, working through the night so that the
image need not be withdrawn from cult. The counter-frame, the acrylic
case of 1976, and nearly fifty studded nails tensioning the cloth
were removed; a consolidant was applied; and the reverse of the
painting was seen for the third time in its documented history.

The findings, as summarized in the peer-reviewed art-historical
literature from the published reports, were these. There is no
varnish. Traces of brush drag are hard to find, which the report
attributed to the \term{aguazo} technique used on \term{sargas}, in
which the cloth is dampened to fix the colour---the same technique
that would explain the bleeding through to the reverse. Two further
tempera media are present: a glue tempera in the sunburst and a resin
tempera in the hands and face, the most lustrous part. The garments
are worked in a technique of glue and ground gold. The palette is
restricted; the black is soot from burnt maize-cob; the blue, often
taken for Maya blue, is achieved with copper carbonates or azurite;
the vermilion of the tunic is from sulphur and mercury, with cochineal
carmine. The support carries a white, milky, irregular
\term{imprimatura}, and its fibres are bast fibres with the
characteristics of hemp---which, the same account observes, closed
the long argument between the partisans of maguey and of palm. Under
``materials foreign to the work'' the report listed various splashes,
including three of wax.\endnote{Cuadriello, ``Cifra, signo y
artilugio'' (same witness), section ``Bajo el lente del
microscopio,'' summarizing and quoting Jos\'e Sol Rosales,
``Dictamen de conservaci\'on de la imagen de la Virgen de
Guadalupe,'' and Jorge Guadarrama Guevara, ``Informe sobre aspectos
t\'ecnicos de conservaci\'on [de la imagen de la Virgen de
Guadalupe],'' both printed in Fernando Ojeda Llanes, \work{La tilma
guadalupana revela sus secretos} (Mexico: Miguel \'Angel Porr\'ua,
2005), pp.~212--222 and 230--238. The two reports themselves were not
consulted; everything above is at the ceiling of Cuadriello's use of
them. \emph{Recorded correction of a common description}: this
examination is often described in devotional and journalistic
literature as a ``secret study'' that was suppressed. Its principal
documents were printed in a book of 2005 and are cited in the
peer-reviewed literature; what was avoided in 1982, on Cuadriello's
account, was public ``alboroto'' during the work itself.}

\subsection{Claims that could not be traced}

The following claims circulate widely. Each was searched for; the
right-hand column records what this study actually found.\endnote{Every row of this table records the result of searches
performed on 2026-07-25 and is bounded by them: a literal search
cannot exclude a publication in an unsearched language, an offline
archive, or a source not indexed by the tools used. ``Not traced''
means this study could not reach the source, not that no source
exists. The Callahan volume's bibliographic identity is reported from
secondary literature and was not verified against a library record;
its contents were not read, and the summary of its conclusions in the
table is reported, not endorsed. On the survival of the cloth,
Garc\'ia Icazbalceta's reply is at \work{Carta}, \S56: the physicians'
1666 opinion that the preservation was miraculous is met with the
observation that far older and more fragile papers survive, and with
the 1795 canons' own statement that ``los colores se han
amortiguado, deslustrado, y en una u otra parte saltado el oro, y el
lienzo sagrado no poco lastimado.''}

\begin{claimtable}
NASA scientists examined the image and declared it inexplicable, or
``living'' & No statement by the United States National Aeronautics
and Space Administration concerning the image was located & Not
traced. No NASA document, release, or attributed statement was
reached. The claim is not repeated here in any form\\
Richard Kuhn, Nobel laureate in chemistry, analysed the fibres in
1936 and found no known pigments & No publication by Kuhn on the
subject was located & Not traced. No published report, laboratory
record, or citation to a primary publication was reached\\
The cloth would have decayed within decades and its survival is
therefore miraculous & Bartolache's 1787 controlled experiment with
woven ayates is the earliest systematic test located & The historical
argument is old and was already answered in 1883 by the observation
that far more fragile papers survive far longer; no modern
materials-science study was located\\
Infrared photography in 1979 showed that the original figure has no
underdrawing and no brushwork & Philip Serna Callahan, \work{The
Tilma under Infra-Red Radiation}, CARA Studies on Popular Devotion,
vol.~II, Guadalupan Studies no.~3 (Washington, D.C., 1981) & Publication
identified; \emph{not read}. Its reported conclusions distinguish an
``original'' figure from later additions (moon, tassel, gold rays,
angel, background), which is a claim of later human overpainting\\
Human figures are reflected in the eyes of the image & Studies by
J.~Aste Tonsmann and earlier observations from 1929 and 1951 are
referred to in the popular literature & Not traced to any
peer-reviewed publication reachable for this study\\
The stars on the mantle map the sky over Mexico on 12 December 1531 &
No published astronomical study was located & Not traced\\
The name ``Guadalupe'' renders a Nahuatl \term{coatlaxopeuh}, ``she
who crushes the serpent'' & No Nahuatl-philological support was
located; the 1649 tract itself prints \term{Guadalupe} & Not traced.
The 1556 testimony has the Spanish residents of Mexico City giving
the hermitage that title (\S3.2)\\
\end{claimtable}

Two of these rows deserve a further sentence, because they are
sometimes cited in Catholic apologetics as though they were settled.
The Callahan study, whichever way its conclusions are finally
assessed, is not a finding that the image is untouched: on the
reported summary it explicitly separates an original figure from
later human additions, which is the same structure of claim as the
1982 conservation report. And the \term{coatlaxopeuh} etymology has
against it the plain fact that the earliest datable testimony about
the name---the 1556 inquiry---attributes it to the Spanish residents
of the city, and that the 1649 Nahuatl tract prints the Spanish
form \term{Guadalupe} without gloss.

\subsection{What material findings can and cannot do}

Suppose every material question were settled tomorrow in the sense
most favourable to the tradition. It would still not be a
demonstration that the Blessed Virgin Mary appeared at Tepeyac; an
unexplained material fact is an unexplained material fact, and the
inference from ``science cannot presently account for this'' to
``therefore God did it in this way'' is not a valid one and has never
been the Church's method. Conversely, suppose every material question
were settled in the sense least favourable. A sixteenth-century
painting on hemp, made by a named or unnamed Indigenous painter,
before which millions have prayed and through which the Gospel
reached a devastated people, would remain one of the most
consequential objects in the history of Christian evangelization, and
the Church's veneration of it would be untouched, because the Church
venerates images without asserting anything about their manufacture.

That is the reason this study declines to make the material question
carry weight it cannot bear---and the reason a Catholic reader loses
nothing by holding the material question open.
